\section{Archive-to-Study Mapping}
\label{sec:archive}
The compact source package maps Study~A to archive V400, Study~B to V401, Study~C to V402 and Study~D to V403. Identifiers 2400001--2400016, 2410001--2410004, 2420001--2420016 and 2430001--2430004 distinguish their respective demand panels. The selected original learning days were 1780001--1780032. Study~C calibrated only on Study~B's already consumed development days; the 16 independent Study~C days were not used to choose mixture weights.

The physical-start accounting is separate from panel sample size. Study~A used two engineering reproductions plus 80 test paths. Study~B charged 92 starts, including two preserved engineering failures, their two authorized replacements, 64 matched learning collections and 24 development paths. Study~C used three engineering paths, 64 calibration paths and 80 independent paths. Study~D used three engineering paths plus 24 development paths. Reusing a reference or evaluating several arms on one day does not add independent demand units.

\section{Implementation and Evidence Package}
The simulator/controller study was implemented in an isolated derivative of the TransitDuet codebase. Tables here concern the stated fixed-fleet contract rather than outcomes from the original TransitDuet policy or another frequency-based reinforcement-learning project. The fast implementation enumerates seven current actions, each followed by its model continuation, with immutable pending commands and independent trial ledgers. The common-reference candidate and simple limits retain that event kernel.

The machine-readable manifest records each table estimate's source file, contrast, endpoint, original interval field and units. The table generator formats existing means and intervals only; it performs no new resampling or traffic rollout. The source package includes the generator and compact summaries so the numerical manuscript can be rebuilt offline. Full traffic reproduction additionally requires the archived simulator input files, native runtime bundle and frozen supervisor artifacts identified by the repository protocols.

\section{Endpoint and Gate Preservation}
Study~A's budget-service result remains passed, and its harmonic-journey result remains failed. Study~B's explicit-retention development screen failed. Both Study~C joint delay/regularity comparisons failed. Study~D's joint development screen failed. These outcomes concern different hypotheses and are retained separately. In particular, the successful budget-setting result is not erased by the failed pressure-allocation comparison, and the latter is not rescued by the joint-admission candidate's improved regularity.

Public departure delay uses positive lateness against original public releases; readiness uses first dispatch eligibility under the operational plan. Restricted journey exposure includes the prescribed waiting and onboard tail for the common generated population. Actual arrival CV$^2$ excludes terminal-return events. These definitions remain fixed across the comparisons rather than being selected to favor a particular operating point.
